\section{Mạng nơ-ron tích hợp thông tin vật lý}
\label{chap:pinn}
Mạng nơ-ron tích hợp thông tin vật lý (PINN), sử dụng một họ hàm
tham số hóa để xấp xỉ nghiệm phương trình vi phân. Thông tin về phương trình,
điều kiện đầu, điều kiện biên và các quan trắc được đưa vào quá trình xác
định tham số. Trong cách xây dựng bằng phần dư mạnh, bài toán vi phân được
chuyển thành một bài toán bình phương tối thiểu phi tuyến.

Chương này gồm ba nội dung. Mục~\ref{sec:cau-truc-loss} xây dựng hàm mất
mát và phân biệt phần dư tại các điểm lấy mẫu với sai số nghiệm trên toàn
miền. Mục~\ref{sec:gradient-pathology} phân tích hình học của hàm mất mát
trong không gian tham số, sự mất cân bằng giữa các građien thành phần và
cơ chế điều chỉnh trọng số. Mục~\ref{sec:spectral-bias} trình bày thiên
kiến phổ, đồng thời phân biệt khả năng biểu diễn của mạng với tốc độ học
các thành phần tần số. Ba cơ sở tài liệu chính lần lượt là công trình của
Raissi và cộng sự~\cite{raissi2019}, Wang và cộng sự~\cite{wang2021},
Rahaman và cộng sự~\cite{rahaman2019}. Các mệnh đề có chứng minh trong
chương được phát biểu với giả thiết riêng, không được hiểu là các bảo đảm
tổng quát cho mọi bài toán PINN.

\section{Cấu trúc hàm mất mát trong PINN}
\label{sec:cau-truc-loss}

\subsection{Bài toán vi phân và họ hàm xấp xỉ}

Cho miền bị chặn $\Omega\subset\mathbb{R}^{d}$ có biên đủ chính quy và
$T>0$. Đặt $Q=\Omega\times(0,T)$,
$\Sigma=\partial\Omega\times(0,T)$. Xét bài toán
\begin{align}
 \partial_tu(\mathbf{x},t)+\mathcal{N}[u;\boldsymbol{\zeta}](\mathbf{x},t)
 &=f(\mathbf{x},t), &&(\mathbf{x},t)\in Q,\label{p4:eq:pde}\\
 u(\mathbf{x},0)&=u_0(\mathbf{x}),&&\mathbf{x}\in\Omega,\label{p4:eq:ic}\\
 \mathcal{B}[u](\mathbf{x},t)&=g(\mathbf{x},t),&&
                   (\mathbf{x},t)\in\Sigma.\label{p4:eq:bc}
\end{align}
Trong đó, $u$ là nghiệm cần tìm; $\mathbf{x}$ và $t$ lần lượt là tọa độ
không gian và thời gian vật lý; $d$ là số chiều không gian;
$\mathcal{N}$ là toán tử vi phân theo không gian, có thể phi tuyến;
$\boldsymbol{\zeta}$ là vectơ hệ số vật lý; $f$ là nguồn đã biết;
$u_0$ là dữ kiện đầu; $\mathcal{B}$ là toán tử biên và $g$ là dữ kiện biên.
Ví dụ, $\mathcal{B}[u]=u$ cho điều kiện Dirichlet, còn
$\mathcal{B}[u]=\nabla_{\mathbf{x}}u\cdot\mathbf{n}$ cho điều kiện Neumann,
với $\mathbf{n}$ là pháp tuyến đơn vị hướng ra ngoài.

Ta giả thiết bài toán đang xét có nghiệm trong một không gian hàm thích
hợp. Tính duy nhất và ổn định phải được xác lập theo từng phương trình
và từng chuẩn sai số; việc xây dựng một hàm mất mát không tự nó chứng minh
các tính chất đó.

Xấp xỉ $u$ bằng mạng truyền thẳng $u_{\boldsymbol{\theta}}$ với đầu vào
$\mathbf{z}=(\mathbf{x},t)\in\mathbb{R}^{d+1}$. Một quy ước cụ thể là
\begin{align}
 \mathbf{a}^{[0]}&=\mathbf{z},\notag\\
 \mathbf{a}^{[l]}&=\sigma\!\left(\mathbf{W}^{[l]}\mathbf{a}^{[l-1]}
                                     +\mathbf{b}^{[l]}\right),
                       &&l=1,\ldots,D-1,\notag\\
 u_{\boldsymbol{\theta}}(\mathbf{z})
 &=\mathbf{W}^{[D]}\mathbf{a}^{[D-1]}+b^{[D]}.
 \label{p4:eq:network}
\end{align}
Trong đó, $D$ là số lớp có tham số, lớp cuối có một đầu ra;
$\mathbf{W}^{[l]}\in\mathbb{R}^{n_l\times n_{l-1}}$ và
$\mathbf{b}^{[l]}\in\mathbb{R}^{n_l}$ là trọng số và độ lệch;
$n_0=d+1$, $n_D=1$; $\sigma$ tác động lên từng thành phần.
Vectơ $\boldsymbol{\theta}\in\mathbb{R}^{P}$ ghép tất cả các tham số,
với $P=\sum_{l=1}^{D}n_l(n_{l-1}+1)$. Chỉ số lớp luôn viết trong ngoặc
vuông; chỉ số $k$ ở $\boldsymbol{\theta}_k$ về sau chỉ bước tối ưu.

\subsection{Phần dư và vai trò của vi phân tự động}

\begin{definition}[Phần dư mạnh]
Giả sử mạng có đủ các đạo hàm xuất hiện trong \eqref{p4:eq:pde}.
Phần dư của phương trình tại bộ tham số $\boldsymbol{\theta}$ là
\begin{equation}
 r_{\boldsymbol{\theta},\boldsymbol{\zeta}}(\mathbf{x},t)
 :=\partial_tu_{\boldsymbol{\theta}}(\mathbf{x},t)
 +\mathcal{N}[u_{\boldsymbol{\theta}};\boldsymbol{\zeta}](\mathbf{x},t)
 -f(\mathbf{x},t).
 \label{p4:eq:residual}
\end{equation}
\end{definition}
Khi hệ số vật lý đã biết, ta viết gọn $r_{\boldsymbol{\theta}}$.
Phần dư bằng không có nghĩa mạng thỏa phương trình tại điểm đang xét;
điều này chưa bao gồm điều kiện đầu và điều kiện biên.

Chẳng hạn, phương trình Burgers có độ nhớt $\nu>0$ và nguồn bằng không
cho
\begin{equation}
 r_{\boldsymbol{\theta}}
 =\partial_tu_{\boldsymbol{\theta}}
 +u_{\boldsymbol{\theta}}\partial_xu_{\boldsymbol{\theta}}
 -\nu\partial_{xx}u_{\boldsymbol{\theta}}.
 \label{p4:eq:burgers-residual}
\end{equation}
Ba số hạng lần lượt biểu diễn biến thiên theo thời gian, vận chuyển phi
tuyến và khuếch tán. Theo cách xây dựng của Raissi và cộng
sự~\cite{raissi2019,raissi2017a}, các đạo hàm này được tính từ cùng mạng
$u_{\boldsymbol{\theta}}$ bằng quy tắc dây chuyền trên đồ thị tính toán.
Phần dư không phải một mạng độc lập có bộ trọng số riêng.

Cần phân biệt đạo hàm theo đầu vào với đạo hàm theo tham số. Để tính
phần dư, ta dùng $\partial_xu_{\boldsymbol{\theta}}$ và
$\partial_{xx}u_{\boldsymbol{\theta}}$; để huấn luyện, ta còn cần các
đạo hàm như $\partial_{\theta_p}\partial_{xx}u_{\boldsymbol{\theta}}$,
trong đó $\theta_p$ là tham số thứ $p$. Vi phân tự động không sử dụng
bước sai phân, nhưng vẫn chịu sai số làm tròn và sai số do điều kiện
số của đồ thị tính toán. Chi phí đạo hàm bậc cao và đạo hàm hỗn hợp
cũng cần được tính vào chi phí huấn luyện.

Với phương trình bậc $m$ theo không gian, một điều kiện đủ thuận tiện
cho phần dư cổ điển là chọn các hàm kích hoạt thuộc $C^m$; dùng
$\tanh$ còn bảo đảm tính trơn ở mọi bậc hữu hạn. Điều kiện này không
phải một định lý rằng mọi đạo hàm bậc cao đều khác không. Ngược lại,
mạng dùng hàm tuyến tính chỉnh lưu ReLU, $\sigma(z)=\max\{0,z\}$, là hàm afin liên tục từng
khúc, nên đạo hàm không gian bậc hai bằng không bên trong mỗi miền
tuyến tính. Với phần dư Poisson thuần túy $-u_{\boldsymbol{\theta}}''-f$,
điều này có thể làm mất thông tin về độ cong tại các điểm phối trí nằm
ngoài điểm gãy. Tuy nhiên, các số hạng điều kiện biên, dữ liệu hoặc
đạo hàm bậc nhất vẫn có thể có građien tham số khác không. Do đó không
thể kết luận toàn bộ hàm mất mát của mọi PINN dùng ReLU đều là hằng số.

\subsection{Điểm phối trí và chiến lược lấy mẫu}

Chọn các tập hữu hạn
\begin{align}
 \mathcal{T}_r&=\{(\mathbf{x}_r^i,t_r^i)\}_{i=1}^{N_r}\subset Q,\notag\\
 \mathcal{T}_0&=\{\mathbf{x}_0^i\}_{i=1}^{N_0}\subset\Omega,\notag\\
 \mathcal{T}_b&=\{(\mathbf{x}_b^i,t_b^i)\}_{i=1}^{N_b}\subset\Sigma.
 \label{p4:eq:sets}
\end{align}
Trong đó, $N_r,N_0,N_b$ là số điểm tương ứng. Các điểm trong
$\mathcal{T}_r$ gọi là điểm phối trí: không cần biết nghiệm đúng tại
đó, vì giá trị đích của phần dư là không. Tại $\mathcal{T}_0$ và
$\mathcal{T}_b$, dữ kiện được cung cấp bởi $u_0$ và $g$.

Trên miền hộp đã chuẩn hóa thành $[0,1]^{d+1}$, lấy mẫu siêu khối Latin
có thể được mô tả bởi
\begin{equation}
 z_{ij}=\frac{\pi_j(i)-1+U_{ij}}{N_r},
 \qquad i=1,\ldots,N_r,\quad j=1,\ldots,d+1.
 \label{p4:eq:lhs}
\end{equation}
Ở đây, $\pi_j$ là một hoán vị ngẫu nhiên của $\{1,\ldots,N_r\}$;
$U_{ij}$ phân bố đều trên $(0,1)$; các hoán vị theo chiều và các biến
$U_{ij}$ được chọn độc lập. Mỗi khoảng phân tầng trên một trục chứa
đúng một tọa độ mẫu. Đây là ý nghĩa chính xác của sự phân bố đều theo
từng chiều, không phải bảo đảm rằng mọi miền con nhiều chiều đều chứa
điểm. Raissi và cộng sự dùng cách lấy mẫu này trong thí nghiệm
Burgers~\cite{raissi2017a}.

Các điểm trong một thiết kế Latin không độc lập với nhau. Vì vậy,
không áp dụng trực tiếp công thức phương sai dành cho các điểm độc lập
cho thiết kế này. Lấy mẫu ngẫu nhiên độc lập, lấy mẫu phân tầng và các
dãy tựa ngẫu nhiên cũng không có hiệu quả tương đương trên mọi hàm:
hiệu quả phụ thuộc tính trơn, số chiều và sự tập trung của phần dư.
Với lớp biên mỏng, cần kiểm tra số điểm thực sự nằm trong lớp, thay vì
chỉ dựa vào tổng số điểm trên toàn miền.

\subsection{Hàm mất mát rời rạc và việc chọn trọng số}

Định nghĩa các thành phần
\begin{align}
 J_{r,N_r}(\boldsymbol{\theta})
 &=\frac{1}{N_r}\sum_{i=1}^{N_r}
       |r_{\boldsymbol{\theta}}(\mathbf{x}_r^i,t_r^i)|^2,
       \label{p4:eq:loss-r}\\
 J_{0,N_0}(\boldsymbol{\theta})
 &=\frac{1}{N_0}\sum_{i=1}^{N_0}
       |u_{\boldsymbol{\theta}}(\mathbf{x}_0^i,0)-u_0(\mathbf{x}_0^i)|^2,
       \label{p4:eq:loss-0}\\
 J_{b,N_b}(\boldsymbol{\theta})
 &=\frac{1}{N_b}\sum_{i=1}^{N_b}
       |\mathcal{B}[u_{\boldsymbol{\theta}}](\mathbf{x}_b^i,t_b^i)
                                  -g(\mathbf{x}_b^i,t_b^i)|^2.
       \label{p4:eq:loss-b}
\end{align}
Trong đó, trị tuyệt đối áp dụng cho phương trình vô hướng; với hệ
phương trình có thể thay bằng chuẩn Euclid và chỉ rõ cách cân bằng
các phương trình. Chỉ số $N_r,N_0,N_b$ nhấn mạnh sự phụ thuộc vào tập
lấy mẫu. Ta viết $J_r,J_0,J_b$ khi không có nguy cơ nhầm lẫn.

Hàm mất mát tổng hợp là
\begin{equation}
 J_N(\boldsymbol{\theta})=
 \lambda_rJ_{r,N_r}(\boldsymbol{\theta})+
 \lambda_0J_{0,N_0}(\boldsymbol{\theta})+
 \lambda_bJ_{b,N_b}(\boldsymbol{\theta}),
 \qquad \lambda_r,\lambda_0,\lambda_b>0.
 \label{p4:eq:loss-total}
\end{equation}
Ký hiệu $N$ đại diện cho bộ kích thước lấy mẫu; các $\lambda$ là trọng
số. Phép chia cho số điểm biến tổng thành trung bình, tránh tăng trọng
số của một thành phần chỉ vì tăng số điểm của nó. Tuy vậy, phép chia
này không tự động cân bằng đơn vị vật lý hoặc độ lớn građien.

Ví dụ, với phép đổi biến $x=\ell\widehat x$, $t=\tau\widehat t$,
$u=U\widehat u$, phương trình Burgers trở thành
\begin{equation}
 \partial_{\widehat t}\widehat u+
 \frac{U\tau}{\ell}\widehat u\,\partial_{\widehat x}\widehat u
 -\frac{\nu\tau}{\ell^2}\partial_{\widehat x\widehat x}\widehat u=0.
 \label{p4:eq:nondimensional}
\end{equation}
Trong đó, $\ell,\tau,U>0$ là các thang không gian, thời gian và nghiệm.
Công thức suy ra bằng cách thay lần lượt
$u_t=(U/\tau)\widehat u_{\widehat t}$,
$u_x=(U/\ell)\widehat u_{\widehat x}$,
$u_{xx}=(U/\ell^2)\widehat u_{\widehat x\widehat x}$ rồi chia cho
$U/\tau$. Vì việc đổi đơn vị có thể thay đổi mất mát phần dư theo bình
phương hệ số tỉ lệ, nên cần chọn thang đo hoặc khử thứ nguyên trước
khi giải thích các trọng số là mức ưu tiên của từng ràng buộc.

Mục tiêu tính toán là tìm $\widehat{\boldsymbol{\theta}}$ sao cho
\begin{equation}
 J_N(\widehat{\boldsymbol{\theta}})
 \leq\inf_{\boldsymbol{\theta}\in\Theta}J_N(\boldsymbol{\theta})
                  +\varepsilon_{\mathrm{opt}},
 \qquad \varepsilon_{\mathrm{opt}}\geq0.
 \label{p4:eq:approx-opt}
\end{equation}
Ở đây $\Theta\subseteq\mathbb{R}^{P}$ là tập tham số được xét và
$\varepsilon_{\mathrm{opt}}$ là sai số tối ưu. Ký hiệu cận dưới đúng không
đòi hỏi cực tiểu phải đạt được. Chỉ khi đã có tính đạt cực tiểu mới
viết $\boldsymbol{\theta}^{*}\in\arg\min J_N$.
Trong thực hành thường không biết $\inf J_N$, nên số bước tối ưu hoặc
chuẩn građien nhỏ chưa đủ để xác nhận một giá trị
$\varepsilon_{\mathrm{opt}}$ cụ thể trong \eqref{p4:eq:approx-opt}.

\subsection{Từ tổng hữu hạn đến chuẩn phần dư trên toàn miền}

Gọi $\mu_r$ là phân phối xác suất dùng làm chuẩn đánh giá trên $Q$.
Hàm mất mát phần dư liên tục là
\begin{equation}
 \mathcal{J}_r(\boldsymbol{\theta})
 =\int_Q|r_{\boldsymbol{\theta}}(\mathbf{z})|^2\,\mathrm{d}\mu_r(\mathbf{z}).
 \label{p4:eq:population-r}
\end{equation}
Với phân phối đều, $\mathrm{d}\mu_r=\mathrm{d}\mathbf{x}\,\mathrm{d}t/(|\Omega|T)$,
nên $\mathcal{J}_r=\|r_{\boldsymbol{\theta}}\|_{L^2(Q)}^2/(|\Omega|T)$.
Tương tự, định nghĩa $\mathcal{J}_0,\mathcal{J}_b$ bằng tích phân theo
các phân phối trên lát đầu và biên, rồi đặt
$\mathcal{J}=\lambda_r\mathcal{J}_r+\lambda_0\mathcal{J}_0+
\lambda_b\mathcal{J}_b$. Với $d=1$, phân phối trên biên có thể chọn đều
trên hai đầu mút và đều theo thời gian.

Nếu $\mathbf{Z}_1,\ldots,\mathbf{Z}_{N_r}$ độc lập, cùng phân phối
$\mu_r$, và $\boldsymbol{\theta}$ cố định, tính tuyến tính của kỳ vọng cho
\begin{equation}
 \mathbb{E}[J_{r,N_r}(\boldsymbol{\theta})]
 =\frac{1}{N_r}\sum_{i=1}^{N_r}
       \mathbb{E}|r_{\boldsymbol{\theta}}(\mathbf{Z}_i)|^2
 =\mathcal{J}_r(\boldsymbol{\theta}).
 \label{p4:eq:unbiased}
\end{equation}
Nếu phần dư thuộc $L^2(\mu_r)$, luật số lớn mạnh còn cho
$J_{r,N_r}(\boldsymbol{\theta})\to\mathcal{J}_r(\boldsymbol{\theta})$
gần chắc chắn khi $N_r\to\infty$. Nếu thêm
$r_{\boldsymbol{\theta}}\in L^4(\mu_r)$, tính độc lập cho
\begin{equation}
 \operatorname{Var}[J_{r,N_r}(\boldsymbol{\theta})]
 =\frac{\operatorname{Var}(|r_{\boldsymbol{\theta}}(\mathbf{Z}_1)|^2)}{N_r}.
 \label{p4:eq:mc-variance}
\end{equation}
Trong đó, $\mathbb{E}$ và $\operatorname{Var}$ lần lượt là kỳ vọng và
phương sai theo phép lấy mẫu. Điều kiện $L^4$ bảo đảm phương sai ở tử
số hữu hạn; chỉ giả thiết $L^2$ chưa đủ cho kết luận phương sai này.

Khi $\widehat{\boldsymbol{\theta}}$ được tối ưu trên chính tập phối trí,
không thể thay nó vào \eqref{p4:eq:unbiased} mà không xét sự phụ thuộc
giữa tham số và mẫu. Sự phụ thuộc ấy có thể gây đánh giá quá lạc quan,
nhưng không chứng minh rằng từng thành phần mất mát luôn bị chệch
xuống trong mọi thuật toán. Một tập kiểm tra mới, độc lập với quá trình
chọn mạng, cho đẳng thức có điều kiện
\begin{equation}
 \mathbb{E}\!\left[J^{\mathrm{test}}_{r,M}
  (\widehat{\boldsymbol{\theta}})\mid\widehat{\boldsymbol{\theta}}\right]
 =\mathcal{J}_r(\widehat{\boldsymbol{\theta}}).
 \label{p4:eq:test-unbiased}
\end{equation}
Ở đây $M$ là số điểm kiểm tra. Tập được dùng để lựa chọn mô hình nhiều
lần không còn là tập kiểm tra độc lập cuối cùng.

Nếu lấy mẫu thích nghi với mật độ $q(\mathbf{z})>0$ so với
$\mu_r$, thì
\begin{equation}
 \mathbb{E}_{\mathbf{Z}\sim q\mu_r}
 \left[\frac{|r_{\boldsymbol{\theta}}(\mathbf{Z})|^2}{q(\mathbf{Z})}\right]
 =\int_Q\frac{|r_{\boldsymbol{\theta}}(\mathbf{z})|^2}{q(\mathbf{z})}
              q(\mathbf{z})\,\mathrm{d}\mu_r
 =\mathcal{J}_r(\boldsymbol{\theta}).
 \label{p4:eq:importance}
\end{equation}
Mật độ $q$ phải tích phân bằng một. Nếu bỏ hệ số $1/q$, ta đang xấp
xỉ một chuẩn có trọng số khác. Với $q$ phụ thuộc lịch sử huấn luyện,
đẳng thức được hiểu có điều kiện theo lịch sử trước khi lấy mẫu mới.

\subsection{Phân rã sai số xấp xỉ, lấy mẫu và tối ưu}

Giả sử trên cùng tập $\Theta$ ta có một cận hữu hạn
\begin{equation}
 \varepsilon_{\mathrm{samp}}
 :=\sup_{\boldsymbol{\theta}\in\Theta}
       |\mathcal{J}(\boldsymbol{\theta})-J_N(\boldsymbol{\theta})|.
 \label{p4:eq:uniform-gap}
\end{equation}
Đây là sai lệch đều giữa mất mát liên tục và rời rạc. Sự hội tụ tại
mỗi tham số cố định trong luật số lớn không tự động cho cận đều này.

\begin{proposition}[Cận cho mất mát liên tục]
Nếu \eqref{p4:eq:approx-opt} và \eqref{p4:eq:uniform-gap} đúng thì
\begin{equation}
 \mathcal{J}(\widehat{\boldsymbol{\theta}})
 \leq\varepsilon_{\mathrm{app}}+
          2\varepsilon_{\mathrm{samp}}+\varepsilon_{\mathrm{opt}},
 \qquad
 \varepsilon_{\mathrm{app}}:=\inf_{\boldsymbol{\theta}\in\Theta}
                                      \mathcal{J}(\boldsymbol{\theta}).
 \label{p4:eq:decomposition}
\end{equation}
\end{proposition}
Trong đó, $\varepsilon_{\mathrm{app}}$ đo khả năng thỏa bài toán bằng
họ hàm đã chọn, khi nghiệm đúng có mất mát liên tục bằng không.

\begin{proof}
Với mọi $\boldsymbol{\theta}\in\Theta$,
\begin{align*}
 \mathcal{J}(\widehat{\boldsymbol{\theta}})
 &\leq J_N(\widehat{\boldsymbol{\theta}})+\varepsilon_{\mathrm{samp}}\\
 &\leq J_N(\boldsymbol{\theta})+
                         \varepsilon_{\mathrm{opt}}+\varepsilon_{\mathrm{samp}}\\
 &\leq\mathcal{J}(\boldsymbol{\theta})+
                         \varepsilon_{\mathrm{opt}}+2\varepsilon_{\mathrm{samp}}.
\end{align*}
Lấy cận dưới đúng theo $\boldsymbol{\theta}$ ở vế phải cho kết quả.
\end{proof}

Một định lý xấp xỉ phổ dụng thường cho phép tăng độ rộng, độ sâu hoặc
thay đổi lớp tham số; nó không chứng minh
$\varepsilon_{\mathrm{app}}=0$ cho một kiến trúc hữu hạn cố định.
Hơn nữa, xấp xỉ tốt chính hàm $u$ theo chuẩn đều chưa tự động cho xấp
xỉ tốt các đạo hàm xuất hiện trong phần dư. Cần chọn chuẩn xấp xỉ và
giả thiết phù hợp với toán tử vi phân.

\subsection{Chặn ổn định: từ phần dư đến sai số nghiệm}

Đặt $e=v-u$, với $v$ là một hàm xấp xỉ đủ trơn. Muốn chuyển
\eqref{p4:eq:decomposition} thành cận sai số nghiệm, cần một bất đẳng
thức ổn định. Ta chứng minh tường minh cho bài toán Poisson một chiều.

\begin{proposition}[Poisson với sai số biên]
Cho $u\in H^2(0,1)$ thỏa $-u''=f$, $u(0)=g_0$, $u(1)=g_1$ và
$v\in H^2(0,1)$. Đặt
\[
 e=v-u,\quad r=-v''-f,\quad
 \beta_0=v(0)-g_0,\quad\beta_1=v(1)-g_1.
\]
Khi đó
\begin{equation}
 \|e\|_{L^2(0,1)}
 \leq\frac{1}{\pi^2}\|r\|_{L^2(0,1)}
      +\sqrt{\frac{\beta_0^2+\beta_1^2}{2}}.
 \label{p4:eq:poisson-bound}
\end{equation}
Nếu $\beta_0=\beta_1=0$ thì còn có
\begin{equation}
 \|e'\|_{L^2(0,1)}\leq\frac{1}{\pi}\|r\|_{L^2(0,1)}.
 \label{p4:eq:poisson-h1}
\end{equation}
\end{proposition}
Ở đây $H^2(0,1)$ là không gian các hàm có đạo hàm yếu đến bậc hai
thuộc $L^2$; $H^1_0(0,1)$ là không gian $H^1$ có vết bằng không tại
hai đầu mút. Dấu phẩy trên hàm là đạo hàm theo $x$.

\begin{proof}
Đặt $\ell(x)=(1-x)\beta_0+x\beta_1$ và $w=e-\ell$.
Khi đó $w\in H^2\cap H^1_0$, $\ell''=0$ và $-w''=r$.
Nhân với $w$ và tích phân từng phần:
\[
 \|w'\|_{L^2}^2
 =\int_0^1rw\,\mathrm{d}x
 \leq\|r\|_{L^2}\|w\|_{L^2}
 \leq\frac{1}{\pi}\|r\|_{L^2}\|w'\|_{L^2}.
\]
Bước cuối là bất đẳng thức Poincaré trên $(0,1)$ với vết bằng không.
Nếu $\|w'\|=0$ thì các cận là hiển nhiên; nếu khác không, chia hai vế
cho $\|w'\|$ cho $\|w'\|\leq\pi^{-1}\|r\|$, rồi áp dụng Poincaré
lần nữa cho $\|w\|\leq\pi^{-2}\|r\|$. Mặt khác,
\[
 \|\ell\|_{L^2}^2
 =\frac{\beta_0^2+\beta_0\beta_1+\beta_1^2}{3}
 \leq\frac{\beta_0^2+\beta_1^2}{2},
\]
vì $2\beta_0\beta_1\leq\beta_0^2+\beta_1^2$.
Bất đẳng thức tam giác $\|e\|\leq\|w\|+\|\ell\|$ cho
\eqref{p4:eq:poisson-bound}. Với sai số biên bằng không thì $w=e$,
nên có thêm \eqref{p4:eq:poisson-h1}.
\end{proof}

Với bài toán này, đặt
$\mathcal{J}=\lambda_r\|r\|_{L^2}^2+
\lambda_b(\beta_0^2+\beta_1^2)/2$. Cauchy--Schwarz trong
$\mathbb{R}^2$ áp dụng cho \eqref{p4:eq:poisson-bound} cho
\begin{equation}
 \|v-u\|_{L^2}
 \leq C_{\mathrm{stab}}\sqrt{\mathcal{J}},
 \qquad
 C_{\mathrm{stab}}=
       \sqrt{\frac{1}{\pi^4\lambda_r}+\frac{1}{\lambda_b}}.
 \label{p4:eq:stability-loss}
\end{equation}
Do đó, nếu các giả thiết của phép phân rã sai số đúng, ta có
\begin{equation}
 \|u_{\widehat{\boldsymbol{\theta}}}-u\|_{L^2}
 \leq C_{\mathrm{stab}}
       \sqrt{\varepsilon_{\mathrm{app}}+
                       2\varepsilon_{\mathrm{samp}}+\varepsilon_{\mathrm{opt}}}.
 \label{p4:eq:total-error}
\end{equation}
Đây là một cận có điều kiện cho mô hình Poisson vừa xét. Không thay
chuẩn phần dư liên tục bằng mất mát trên một tập hữu hạn rồi giữ
nguyên kết luận mà bỏ số hạng lấy mẫu.

\paragraph{Không đồng nhất một cận thô với hằng số ổn định tối ưu.}
Xét $-\nu e''+e'=r$ trên $(0,1)$, $e(0)=e(1)=0$, $\nu>0$.
Nhân với $e$ rồi tích phân cho
$\nu\|e'\|_{L^2}^2=\int re$, vì
$\int e'e=[e^2/2]_0^1=0$. Lập luận Poincaré như trên cho
\begin{equation}
 \|e'\|_{L^2}\leq\frac{\|r\|_{L^2}}{\pi\nu},
 \qquad \|e\|_{L^2}\leq\frac{\|r\|_{L^2}}{\pi^2\nu}.
 \label{p4:eq:coarse-convection}
\end{equation}
Hai bất đẳng thức đúng, nhưng chưa chứng minh chuẩn của toán tử nghịch
đảo theo $L^2$ thực sự tăng như $1/\nu$. Thật vậy, chọn trọng số
$\omega(x)=\exp(-x)$, nhân phương trình với $\omega e$ và tích phân:
\begin{align}
 \int_0^1(-\nu e''+e')\omega e\,\mathrm{d}x
 &=\nu\int_0^1\omega(e')^2\,\mathrm{d}x
    -\frac{\nu}{2}\int_0^1\omega''e^2\,\mathrm{d}x
    -\frac12\int_0^1\omega'e^2\,\mathrm{d}x\notag\\
 &=\nu\int_0^1\omega(e')^2\,\mathrm{d}x
    +\frac{1-\nu}{2}\int_0^1\omega e^2\,\mathrm{d}x
 =\int_0^1\omega re\,\mathrm{d}x.
 \label{p4:eq:weighted-energy}
\end{align}
Các số hạng biên triệt tiêu, còn $\omega'=-\omega$,
$\omega''=\omega$. Đặt $\|h\|_\omega^2=\int_0^1\omega|h|^2$.
Với $0<\nu\leq1/2$, bỏ số hạng không âm và dùng Cauchy--Schwarz:
\[
 \tfrac14\|e\|_\omega^2\leq\|r\|_\omega\|e\|_\omega.
\]
Vì $\exp(-1)\leq\omega\leq1$, suy ra cận đồng đều theo $\nu$:
\begin{equation}
 \boxed{\|e\|_{L^2}\leq4\exp(1/2)\|r\|_{L^2},
                       \qquad 0<\nu\leq1/2.}
 \label{p4:eq:uniform-convection}
\end{equation}
Cận này không kiểm soát đồng đều đạo hàm $e'$, và không loại trừ sự
hình thành lớp biên. Nó cho thấy cần phân biệt sự suy giảm chất lượng
của một phép ước lượng với sự mất ổn định thật sự trong chuẩn đang xét.

\subsection{Ràng buộc cứng và bài toán ngược}

Trong \eqref{p4:eq:loss-total}, điều kiện biên được áp đặt bằng số hạng
phạt nên gọi là ràng buộc mềm. Với điều kiện Dirichlet, nếu dựng được
các hàm đủ trơn $G,\psi$ thỏa $G|_\Sigma=g$, $\psi|_\Sigma=0$, thì
\begin{equation}
 v_{\boldsymbol{\theta}}(\mathbf{x},t)
 =G(\mathbf{x},t)+\psi(\mathbf{x},t)u_{\boldsymbol{\theta}}(\mathbf{x},t)
 \label{p4:eq:hard}
\end{equation}
thỏa biên với mọi tham số: trên biên, số hạng chứa $\psi$ bằng không.
Ở đây $G$ là hàm nâng dữ kiện biên, $\psi$ là nhân tử triệt tiêu trên
biên. Trên $(a,b)$ với biên thuần nhất có thể dùng
$\psi(x)=(x-a)(b-x)$. Cấu trúc này không tự động áp đặt điều kiện
Neumann; với biên và điều kiện đầu đồng thời còn phải xét tính tương
thích. Nó cũng thay đổi đạo hàm của hàm thử, chẳng hạn
$(\psi u_{\boldsymbol{\theta}})''=\psi u_{\boldsymbol{\theta}}''+
2\psi'u_{\boldsymbol{\theta}}'+\psi''u_{\boldsymbol{\theta}}$.
Vì thế, giảm số ràng buộc mềm không đồng nghĩa bảo đảm tối ưu dễ hơn.

Trong bài toán thuận, $\boldsymbol{\zeta}$ đã biết. Trong bài toán
ngược, thêm dữ liệu quan trắc
$\mathcal{T}_u=\{(\mathbf{x}_u^i,t_u^i,y_i)\}_{i=1}^{N_u}$ và đặt
\begin{equation}
 J_u(\boldsymbol{\theta})=\frac{1}{N_u}\sum_{i=1}^{N_u}
 |u_{\boldsymbol{\theta}}(\mathbf{x}_u^i,t_u^i)-y_i|^2.
 \label{p4:eq:inverse-data}
\end{equation}
Trong đó $y_i$ là giá trị quan trắc, có thể có nhiễu. Ta tối ưu đồng
thời $\boldsymbol{\theta}$ và $\boldsymbol{\zeta}$ trong
$\lambda_uJ_u+\lambda_rJ_r+\lambda_0J_0+\lambda_bJ_b$, giữ các thành
phần có dữ kiện tương ứng. Cách ghép dữ liệu với phần dư là nội dung
của mô hình nhận dạng phương trình trong
\cite{raissi2019,raissi2017b}.

Ví dụ, phần dư Burgers có hệ số chưa biết là
\begin{equation}
 r_{\boldsymbol{\theta},\boldsymbol{\zeta}}
 =u_{\boldsymbol{\theta},t}+
 \zeta_1u_{\boldsymbol{\theta}}u_{\boldsymbol{\theta},x}
 -\zeta_2u_{\boldsymbol{\theta},xx}.
 \label{p4:eq:inverse-burgers}
\end{equation}
Hai hệ số $\zeta_1,\zeta_2$ có thể nhận đạo hàm qua cùng đồ thị tính
toán. Tuy nhiên, khả năng tính đạo hàm không bảo đảm nhận dạng duy
nhất: nếu $u\equiv0$, phương trình trên đúng với mọi cặp hệ số.
Tổng quát hơn, sự phân biệt hai hệ số đòi hỏi các thông tin gắn với
$u u_x$ và $-u_{xx}$ không suy biến trên dữ liệu. Thông tin vật lý có
thể hỗ trợ nhận dạng nhưng không tự động biến mọi bài toán ngược
thành bài toán đặt chỉnh.

